\section{Motivation}
\label{sec:motivation}
The rapid proliferation of the Internet of Things (\gls{iot}) and edge computing has spurred significant interest in running machine learning models closer to the source on resource-constrained devices, a paradigm shift driven by the emergence of \gls{tinyml}~\cite{razaghi2026cefedavg, sanchez2024federated}. By deploying compact machine learning architectures directly onto low-power microcontrollers, edge intelligence enables real-time on-device training and inference. This localized execution is highly valuable for distributed edge applications like environmental monitoring, as it drastically reduces the volume of data that must be transmitted over the backhaul network, moving away from high-latency and bandwidth-heavy cloud-centric systems.

However, in many IoT deployments, long-range, battery-powered sensing must rely on Low-Power Wide-Area Networks (\gls{lpwan}) such as \gls{lorawan}~\cite{razaghi2026cefedavg}. Under standard regulations (such as the European 868\,MHz band), \gls{lorawan} operates under extreme physical constraints, including a strict 1\,\% duty cycle limit, highly restricted uplink bandwidth (often only a few hundred bits per second), and long inter-packet intervals. Sending continuous streams of raw environmental measurements over these links is highly inefficient, drains edge device batteries rapidly, and leads to severe network congestion.

In addition to physical network bottlenecks, transmitting large amounts of raw sensor data to a central cloud server introduces major privacy and security vulnerabilities~\cite{sanchez2024federated}. For sensitive industrial, corporate, or residential applications, keeping private and sensitive data close to the machine is essential for complying with modern privacy acts like the GDPR~\cite{ikram2026empirical, razaghi2026cefedavg}. 

To simultaneously address communication constraints and privacy concerns, Federated Learning (\gls{fl}) has emerged as a compelling decentralized alternative~\cite{ikram2026empirical}. Under the principle of model training where the data is not shared but the models are, \gls{fl} allows heterogeneous edge nodes to collaboratively train a global model~\cite{sanchez2024federated}. Devices keep their raw sensor data locally on premises and exchange only small, compressed model weight updates with a central server for aggregation~\cite{ikram2026empirical}. This hybrid approach leverages the collective intelligence of the network while minimizing transmission airtime and preserving data privacy, making it highly suitable for resource-constrained \gls{lorawan} TinyML deployments.
